\section{Related Work and Novelty Boundary}
\subsection{RF Fingerprinting and Receiver Dependence}
WiSig supplies a multi-receiver, multi-day resource for studying receiver and channel variation in WiFi fingerprinting~\cite{wisig}, collected using the ORBIT testbed~\cite{orbit}. Its availability makes receiver holdout possible, but does not make arbitrary subsets of stored packets equivalent to acquisition episodes. The present result uses a support-qualified closed-set subset, not the full transmitter universe in the WiSig paper.

Signal processing and representation choices are central to RF robustness. Shen et al.\ studied LoRa fingerprinting with spectrograms and carrier-frequency-offset effects~\cite{shenlora}; subsequent receiver-agnostic and collaborative work explicitly addresses receiver variation~\cite{shen}. These studies establish that receiver-robust RF identification is not a new problem. We do not reproduce their reported accuracy or rename our source-domain adversarial baseline as their algorithm. The frozen WiSig input is a short I/Q window, whereas the inspected Shen implementation uses channel-independent time--frequency preprocessing and a different network geometry.

\subsection{Source-Domain Generalization}
Adversarial domain learning promotes representations useful for a task while reducing domain discriminability~\cite{dann}. Correlation alignment penalizes differences in feature covariance~\cite{coral}. Group distributionally robust optimization emphasizes worst-group training loss and depends on appropriate regularization~\cite{groupdro}. Our CORAL-style, DANN-style, and GroupDRO baselines apply these objectives using source receiver domains only. This is a protocol-specific implementation of established ideas, not an exhaustive reproduction of all original experimental settings. No unlabeled held-out receiver batch enters these source-only objectives.

\subsection{Test-Time Statistics and Classifier Adaptation}
Tent adapts a model by entropy minimization through normalization-related parameters and statistics~\cite{tent}. Its original implementation contract is not interchangeable with arbitrary gradient updates to a frozen network. The backbone here uses GroupNorm~\cite{groupnorm}, with no BatchNorm state. AdaBN and the verified BatchNorm-based Tent implementation were therefore excluded rather than retrofitted. That choice protects the controlled comparison but leaves a material baseline-coverage limitation.

T3A updates classifier templates using model-derived pseudo-labels and entropy-filtered features~\cite{t3a}. We reuse the frozen bounded-support implementation, resetting its state for each receiver and seed. It receives the same support bank as receiver-context conditioning. We do not introduce a new prototype method or attribute T3A's mechanism to this work.

\subsection{Receiver Calibration and Hardware-Shift Evaluation}
We use receiver adaptation for a model's use of unlabeled held-out receiver information. This should not be confused with instrument calibration, synchronized physical calibration sessions, or calibrated predictive probabilities. Expected calibration error (ECE) describes agreement between confidence and correctness~\cite{calibration}; it answers a different question from transmitter classification accuracy. Hardware-family summaries are descriptive because the evaluated receivers represent only a few families. The novelty boundary is a systematic, information-budget-aware RF evaluation with negative context results, not a claim that receiver adaptation, domain holdout, or prototype adjustment is new.
